Normalise the marble count before turn n by multiplying its residue modulo 15 by the inverse of $2^n$. The next normalised residue is either $8r$ or $8(r+1)$ modulo 15, each with probability one half. Multiplication by eight preserves divisibility by three and five. Set H equal to one on residues divisible by five, including zero, and zero on residues divisible by three but not five. For the eight surviving residues $1,2,4,7,8,11,13,14$, solve $H(r)=\dfrac{[H(8r)+H(8r+8)]}{2}$, with arguments reduced modulo 15. The resulting value at the initial normalised residue $\dfrac{2}{2}=1$ is $\boxed{\dfrac{2}{7}}$. The computation audit includes exact rational elimination of the equivalent turn-phase system. Absorption occurs almost surely, since from every surviving state some bounded-length coin sequence reaches a stopping residue, with a uniformly positive probability.
